% Options for packages loaded elsewhere
\PassOptionsToPackage{unicode}{hyperref}
\PassOptionsToPackage{hyphens}{url}
\documentclass[
  12pt,
]{ctexart}
\usepackage{xcolor}
\usepackage{amsmath,amssymb}
\setcounter{secnumdepth}{-\maxdimen} % remove section numbering
\usepackage{iftex}
\ifPDFTeX
  \usepackage[T1]{fontenc}
  \usepackage[utf8]{inputenc}
  \usepackage{textcomp} % provide euro and other symbols
\else % if luatex or xetex
  \usepackage{unicode-math} % this also loads fontspec
  \defaultfontfeatures{Scale=MatchLowercase}
  \defaultfontfeatures[\rmfamily]{Ligatures=TeX,Scale=1}
\fi
\usepackage{lmodern}
\ifPDFTeX\else
  % xetex/luatex font selection
\fi
% Use upquote if available, for straight quotes in verbatim environments
\IfFileExists{upquote.sty}{\usepackage{upquote}}{}
\IfFileExists{microtype.sty}{% use microtype if available
  \usepackage[]{microtype}
  \UseMicrotypeSet[protrusion]{basicmath} % disable protrusion for tt fonts
}{}
\makeatletter
\@ifundefined{KOMAClassName}{% if non-KOMA class
  \IfFileExists{parskip.sty}{%
    \usepackage{parskip}
  }{% else
    \setlength{\parindent}{0pt}
    \setlength{\parskip}{6pt plus 2pt minus 1pt}}
}{% if KOMA class
  \KOMAoptions{parskip=half}}
\makeatother
\usepackage{color}
\usepackage{fancyvrb}
\newcommand{\VerbBar}{|}
\newcommand{\VERB}{\Verb[commandchars=\\\{\}]}
\DefineVerbatimEnvironment{Highlighting}{Verbatim}{commandchars=\\\{\}}
% Add ',fontsize=\small' for more characters per line
\newenvironment{Shaded}{}{}
\newcommand{\AlertTok}[1]{\textcolor[rgb]{1.00,0.00,0.00}{\textbf{#1}}}
\newcommand{\AnnotationTok}[1]{\textcolor[rgb]{0.38,0.63,0.69}{\textbf{\textit{#1}}}}
\newcommand{\AttributeTok}[1]{\textcolor[rgb]{0.49,0.56,0.16}{#1}}
\newcommand{\BaseNTok}[1]{\textcolor[rgb]{0.25,0.63,0.44}{#1}}
\newcommand{\BuiltInTok}[1]{\textcolor[rgb]{0.00,0.50,0.00}{#1}}
\newcommand{\CharTok}[1]{\textcolor[rgb]{0.25,0.44,0.63}{#1}}
\newcommand{\CommentTok}[1]{\textcolor[rgb]{0.38,0.63,0.69}{\textit{#1}}}
\newcommand{\CommentVarTok}[1]{\textcolor[rgb]{0.38,0.63,0.69}{\textbf{\textit{#1}}}}
\newcommand{\ConstantTok}[1]{\textcolor[rgb]{0.53,0.00,0.00}{#1}}
\newcommand{\ControlFlowTok}[1]{\textcolor[rgb]{0.00,0.44,0.13}{\textbf{#1}}}
\newcommand{\DataTypeTok}[1]{\textcolor[rgb]{0.56,0.13,0.00}{#1}}
\newcommand{\DecValTok}[1]{\textcolor[rgb]{0.25,0.63,0.44}{#1}}
\newcommand{\DocumentationTok}[1]{\textcolor[rgb]{0.73,0.13,0.13}{\textit{#1}}}
\newcommand{\ErrorTok}[1]{\textcolor[rgb]{1.00,0.00,0.00}{\textbf{#1}}}
\newcommand{\ExtensionTok}[1]{#1}
\newcommand{\FloatTok}[1]{\textcolor[rgb]{0.25,0.63,0.44}{#1}}
\newcommand{\FunctionTok}[1]{\textcolor[rgb]{0.02,0.16,0.49}{#1}}
\newcommand{\ImportTok}[1]{\textcolor[rgb]{0.00,0.50,0.00}{\textbf{#1}}}
\newcommand{\InformationTok}[1]{\textcolor[rgb]{0.38,0.63,0.69}{\textbf{\textit{#1}}}}
\newcommand{\KeywordTok}[1]{\textcolor[rgb]{0.00,0.44,0.13}{\textbf{#1}}}
\newcommand{\NormalTok}[1]{#1}
\newcommand{\OperatorTok}[1]{\textcolor[rgb]{0.40,0.40,0.40}{#1}}
\newcommand{\OtherTok}[1]{\textcolor[rgb]{0.00,0.44,0.13}{#1}}
\newcommand{\PreprocessorTok}[1]{\textcolor[rgb]{0.74,0.48,0.00}{#1}}
\newcommand{\RegionMarkerTok}[1]{#1}
\newcommand{\SpecialCharTok}[1]{\textcolor[rgb]{0.25,0.44,0.63}{#1}}
\newcommand{\SpecialStringTok}[1]{\textcolor[rgb]{0.73,0.40,0.53}{#1}}
\newcommand{\StringTok}[1]{\textcolor[rgb]{0.25,0.44,0.63}{#1}}
\newcommand{\VariableTok}[1]{\textcolor[rgb]{0.10,0.09,0.49}{#1}}
\newcommand{\VerbatimStringTok}[1]{\textcolor[rgb]{0.25,0.44,0.63}{#1}}
\newcommand{\WarningTok}[1]{\textcolor[rgb]{0.38,0.63,0.69}{\textbf{\textit{#1}}}}
\usepackage{longtable,booktabs,array}
\usepackage{calc} % for calculating minipage widths
% Correct order of tables after \paragraph or \subparagraph
\usepackage{etoolbox}
\makeatletter
\patchcmd\longtable{\par}{\if@noskipsec\mbox{}\fi\par}{}{}
\makeatother
% Allow footnotes in longtable head/foot
\IfFileExists{footnotehyper.sty}{\usepackage{footnotehyper}}{\usepackage{footnote}}
\makesavenoteenv{longtable}
\setlength{\emergencystretch}{3em} % prevent overfull lines
\providecommand{\tightlist}{%
  \setlength{\itemsep}{0pt}\setlength{\parskip}{0pt}}
\usepackage{bookmark}
\IfFileExists{xurl.sty}{\usepackage{xurl}}{} % add URL line breaks if available
\urlstyle{same}
\hypersetup{
  hidelinks,
  pdfcreator={LaTeX via pandoc}}

\author{SCX}
\date{2026}

\begin{document}

\section{SCX
净室代码检查报告}\label{scx-ux51c0ux5ba4ux4ee3ux7801ux68c0ux67e5ux62a5ux544a}

\begin{quote}
检查日期：2026-06-26 检查工具：\texttt{src/scx/}
及周边目录的自动模式扫描 检查目标：确认 SCX 代码不含学校/课题组特定引用
\end{quote}

\begin{center}\rule{0.5\linewidth}{0.5pt}\end{center}

\subsection{检查范围与模式}\label{ux68c0ux67e5ux8303ux56f4ux4e0eux6a21ux5f0f}

\begin{longtable}[]{@{}
  >{\raggedright\arraybackslash}p{(\linewidth - 2\tabcolsep) * \real{0.5000}}
  >{\raggedright\arraybackslash}p{(\linewidth - 2\tabcolsep) * \real{0.5000}}@{}}
\toprule\noalign{}
\begin{minipage}[b]{\linewidth}\raggedright
扫描目标
\end{minipage} & \begin{minipage}[b]{\linewidth}\raggedright
扫描模式
\end{minipage} \\
\midrule\noalign{}
\endhead
\bottomrule\noalign{}
\endlastfoot
\texttt{src/scx/core/} & \texttt{cxshu}, \texttt{simu\_data},
\texttt{EGP}, \texttt{egp}, \texttt{AlGaN}, \texttt{孝感},
\texttt{Xiaogan}, \texttt{VASP}, \texttt{DFT} \\
\texttt{src/scx/state/} & 同上 \\
\texttt{src/scx/expert/} & 同上 \\
\texttt{src/scx/valuation/} & 同上 \\
\texttt{src/scx/action/} & 同上 \\
\texttt{src/scx/utils/} & 同上 \\
\texttt{tests/} & 同上 \\
\texttt{experiments/} & 同上 \\
\texttt{theory/} & 同上 \\
\end{longtable}

扫描工具：Python 脚本遍历所有 \texttt{.py}, \texttt{.md}, \texttt{.txt}
文件，列级别精确匹配。

\begin{center}\rule{0.5\linewidth}{0.5pt}\end{center}

\subsection{结果总览}\label{ux7ed3ux679cux603bux89c8}

\begin{longtable}[]{@{}lll@{}}
\toprule\noalign{}
目录 & 判定 & 发现数 \\
\midrule\noalign{}
\endhead
\bottomrule\noalign{}
\endlastfoot
\textbf{\texttt{src/scx/}} (核心 Python 包) & \textbf{CLEAN} &
\textbf{0} \\
\textbf{\texttt{tests/}} (测试套件) & \textbf{CLEAN} & \textbf{0} \\
\textbf{\texttt{experiments/}} (实验代码) & \textbf{CLEAN} &
\textbf{0} \\
\texttt{theory/} (理论文档) & 含理论关联引用 & 23 处（详见下文） \\
\end{longtable}

\begin{center}\rule{0.5\linewidth}{0.5pt}\end{center}

\subsection{核心 Python
包：确认洁净}\label{ux6838ux5fc3-python-ux5305ux786eux8ba4ux6d01ux51c0}

\texttt{src/scx/} 下的所有 \texttt{.py}
文件\textbf{未发现}任何以下模式： -
学校服务器路径（\texttt{/home/cxshu/}, \texttt{W:/simu\_data/}） - EGP
项目代码引用 - DFT/VASP 硬编码 - 课题组或学校特定注释/标记

\subsubsection{已检查的源文件}\label{ux5df2ux68c0ux67e5ux7684ux6e90ux6587ux4ef6}

\begin{verbatim}
src/scx/__init__.py
src/scx/core/__init__.py
src/scx/core/config.py
src/scx/core/framework.py
src/scx/core/metrics.py
src/scx/state/__init__.py
src/scx/state/space.py
src/scx/state/discovery.py
src/scx/state/assignment.py
src/scx/state/metrics.py
src/scx/expert/__init__.py
src/scx/expert/registry.py
src/scx/expert/reliability.py
src/scx/expert/router.py
src/scx/expert/conflict.py
src/scx/valuation/__init__.py
src/scx/valuation/learnability.py
src/scx/valuation/noise_score.py
src/scx/valuation/redundancy.py
src/scx/valuation/classifier.py
src/scx/valuation/state_value.py
src/scx/action/__init__.py
src/scx/action/policy.py
src/scx/action/acquisition.py
src/scx/action/compress.py
src/scx/utils/__init__.py
src/scx/utils/helpers.py
src/scx/utils/data_loader.py
src/scx/utils/visualization.py
src/scx/utils/evaluation.py
\end{verbatim}

\textbf{结论：SCX Python 包在净室条件下独立开发，与学校/EGP/DFT
代码无交叉污染。}

\begin{center}\rule{0.5\linewidth}{0.5pt}\end{center}

\subsection{测试与实验：确认洁净}\label{ux6d4bux8bd5ux4e0eux5b9eux9a8cux786eux8ba4ux6d01ux51c0}

\texttt{tests/} 和 \texttt{experiments/}
下的所有文件同样洁净，无学校特定引用。

\begin{center}\rule{0.5\linewidth}{0.5pt}\end{center}

\subsection{理论文档中的引用说明}\label{ux7406ux8bbaux6587ux6863ux4e2dux7684ux5f15ux7528ux8bf4ux660e}

\texttt{theory/} 目录下的数学框架文档中存在 DFT/EGP
引用，但\textbf{这是预期的理论关联}，属于论文线中 Paper 4 (SCX) 与 Paper
1-3 (EGP) 的自然学术衔接，而非代码依赖。

具体文件：

\begin{longtable}[]{@{}
  >{\raggedright\arraybackslash}p{(\linewidth - 4\tabcolsep) * \real{0.2727}}
  >{\raggedright\arraybackslash}p{(\linewidth - 4\tabcolsep) * \real{0.4545}}
  >{\raggedright\arraybackslash}p{(\linewidth - 4\tabcolsep) * \real{0.2727}}@{}}
\toprule\noalign{}
\begin{minipage}[b]{\linewidth}\raggedright
文件
\end{minipage} & \begin{minipage}[b]{\linewidth}\raggedright
引用类型
\end{minipage} & \begin{minipage}[b]{\linewidth}\raggedright
说明
\end{minipage} \\
\midrule\noalign{}
\endhead
\bottomrule\noalign{}
\endlastfoot
\texttt{theory/README.md} & DFT (1 处) &
表格中梳理各方法的数据需求，提及 DFT 作为对比基线 \\
\texttt{theory/propositions/05\_expert\_governance\_protocol.md} & DFT
(21 处), EGP (2 处) & 治理协议的锚定验证步骤讨论了如何用 DFT 作为 ground
truth；这是 Paper 4 与 Paper 1-3 的理论衔接 \\
\end{longtable}

\textbf{这些理论引用不构成代码层面的交叉污染}，它们是学术论文中引用前期工作的正常行为，与净室代码检查的目标无关。

\begin{center}\rule{0.5\linewidth}{0.5pt}\end{center}

\subsection{验证命令}\label{ux9a8cux8bc1ux547dux4ee4}

使用以下命令可复现本检查（需 Python 3.9+）：

\begin{Shaded}
\begin{Highlighting}[]
\ExtensionTok{python} \AttributeTok{{-}c} \StringTok{"}
\StringTok{import os}
\StringTok{suspicious = []}
\StringTok{for root, dirs, files in os.walk(\textquotesingle{}src/scx\textquotesingle{}):}
\StringTok{    for f in files:}
\StringTok{        if f.endswith(\textquotesingle{}.py\textquotesingle{}):}
\StringTok{            path = os.path.join(root, f)}
\StringTok{            with open(path) as fh:}
\StringTok{                content = fh.read()}
\StringTok{            for pattern in [\textquotesingle{}cxshu\textquotesingle{}, \textquotesingle{}simu\_data\textquotesingle{}, \textquotesingle{}EGP\textquotesingle{}, \textquotesingle{}egp\textquotesingle{}, \textquotesingle{}AlGaN\textquotesingle{}, \textquotesingle{}孝感\textquotesingle{}, \textquotesingle{}Xiaogan\textquotesingle{}, \textquotesingle{}VASP\textquotesingle{}, \textquotesingle{}DFT\textquotesingle{}]:}
\StringTok{                if pattern in content:}
\StringTok{                    suspicious.append((path, pattern))}
\StringTok{if suspicious:}
\StringTok{    for p, pat in suspicious:}
\StringTok{        print(f\textquotesingle{}FOUND: \{pat\} in \{p\}\textquotesingle{})}
\StringTok{else:}
\StringTok{    print(\textquotesingle{}CLEAN: No school{-}specific references found\textquotesingle{})}
\StringTok{"}
\end{Highlighting}
\end{Shaded}

\begin{center}\rule{0.5\linewidth}{0.5pt}\end{center}

\subsection{最终判定}\label{ux6700ux7ec8ux5224ux5b9a}

\begin{verbatim}
┌─────────────────────────────────────────────────────┐
│  CLEAN ROOM CHECK: PASSED                           │
│  Core SCX Python package: CLEAN (0 findings)        │
│  Tests:                CLEAN (0 findings)            │
│  Experiments:          CLEAN (0 findings)            │
│  Theory docs:         Contains academic references  │
│                        to EGP/DFT (expected)         │
└─────────────────────────────────────────────────────┘
\end{verbatim}

\end{document}
